% Draft only. Do not include in main.tex until the corresponding result files exist.
\section{Supplementary Evaluation Protocol}
\label{sec:supplementary_protocol}

The supplementary evaluation separates implementation provenance, algorithmic
comparison, and robustness analysis. A complete training run, rather than an
episode, time step, or individual cell, is treated as the independent replicate.
The learned controllers are trained with five prespecified random seeds
($7,17,29,43,71$) for each of the two initial-SOC sets. Runs sharing a seed and
SOC set form a block so that controller comparisons use identical environment
initialization and evaluation conditions.

\subsection{Implementation provenance}

The paper-aligned configuration uses 1200 training episodes and updates the
target networks every 50 training episodes. The reward contains the time,
balancing, voltage, and temperature terms in (9)--(13). Separate diagnostic runs
with the legacy declared configuration (800 episodes and an update interval of
200) and the legacy effective implementation (target copying at every learning
step) are used only to establish the provenance of the original figures. They
are excluded from the confirmatory method comparison.

The provenance runs retain the legacy mixed initialization, which combines a
specified SOC with a voltage-based state and a subsequent negative-electrode
concentration overwrite. All confirmatory runs instead use PyBaMM's
SOC-consistent initialization for both electrodes. Results from these two
initialization modes are not pooled.

Before confirmatory execution, the three PSO-identified parameter dictionaries
must be archived in a machine-readable file and loaded into the corresponding
cell models. The parameter-file hash is stored with every run. If these three
parameter sets cannot be recovered, the evaluation must be relabeled as a
shared-parameter Chen2020 simulation and cannot support a claim of validation
under cell-specific aging variation.

The released observation transformation is documented explicitly for
reproducibility: $SOC_i$ is mapped to $2(SOC_i-0.5)$, voltage is mapped to
$(V_i-3.5)/1$, and temperature is mapped to $(T_i-308)/11$. These constants are
held fixed across all controller comparisons.

\subsection{Matched-controller comparison}

MBA-RL/QMIX is compared with VDN, a matched-architecture DQL controller,
proportional SOC feedback, and per-cell CC--CV. All controllers use the same
identified cell models, independently controllable current-source interface,
action limits, decision interval, state information, safety constraints, and
initial-SOC sets. Two outcomes are reported for every controller: the first time
at which $\sigma\leq\beta$ and the first time at which all cells satisfy
$SOC_i\geq SOC_{ref}$. A run that does not attain an outcome within 50 control
cycles is retained as a failure.

The matched DQL baseline is a centralized single-agent recurrent controller.
It observes the nine-dimensional joint pack state and selects one element from
the Cartesian product of the three per-cell current-action sets. Thus, it uses
the same independently controllable current sources and state-dependent action
masks as QMIX, but it does not use decentralized agent utilities or a mixing
network.

\subsection{Ablation and threshold sensitivity}

The reward ablation removes, one at a time, the time, balancing, or safety reward
term while holding the network, training budget, and environment fixed. The
hard action mask remains active in the no-safety-reward condition. Local
temperature-threshold sensitivity is evaluated at 306, 309, and 312~K. These
experiments assess dependence on the selected reward and threshold; they do not
constitute hardware safety validation.

\subsection{Metrics and statistical reporting}

Primary outcomes are time-to-balance, time-to-reference-SOC, and joint-terminal
success. Secondary outcomes are terminal SOC standard deviation, maximum
voltage, maximum temperature, safety-limit violation counts, and charge/discharge
energy throughput. For learned controllers, each outcome is summarized using
the mean, standard deviation, and a bootstrap 95\% confidence interval across
independent training seeds. Paired differences use runs with the same seed and
SOC set. Energy throughput is reported as a descriptive loading measure and is
not interpreted as converter efficiency.

% RESULT PLACEHOLDER: Insert the completed run counts and provenance statement.
% RESULT PLACEHOLDER: Insert the matched-controller table and confidence intervals.
% RESULT PLACEHOLDER: Insert reward ablation results.
% RESULT PLACEHOLDER: Insert temperature-threshold sensitivity results.
